\documentclass[10pt,a4paper]{article}
\usepackage{xeCJK}
\usepackage{fontspec}
\usepackage{xcolor}
\usepackage{fancyhdr}
\usepackage{booktabs}
\usepackage{array}
\usepackage{amssymb}
\usepackage{geometry}
\usepackage[protrusion=true]{microtype}
\geometry{margin=1.6cm,top=1.8cm,headheight=15pt}
\linespread{1.02}

\setmainfont{Baskerville}
\setCJKmainfont{Songti SC}
\newfontfamily{\starfont}{Songti SC}

\definecolor{wrongred}{RGB}{198,40,40}
\definecolor{rulegray}{RGB}{200,200,200}

\setlength{\emergencystretch}{3em}
\setlength{\fboxsep}{3pt}
\setlength{\parindent}{0pt}

\newcommand{\checkmarkbox}{\fbox{\rule{0pt}{0.8ex}\rule{0.8ex}{0pt}}}
\newcommand{\STAR}{\textcolor{wrongred}{\starfont ★}}
\newcommand{\bk}{\rule[-0.55ex]{1.0em}{0.4pt}}
\newcommand{\A}[1]{\textbf{#1}}
\newcommand{\W}[1]{\textcolor{wrongred}{#1}}
\newcommand{\OK}{\textcolor{wrongred}{\ensuremath{\checkmark}}}
\newcommand{\NO}{\textcolor{wrongred}{$\times$}}
\newcolumntype{L}[1]{>{\raggedright\arraybackslash}p{#1}}

\pagestyle{fancy}
\fancyhf{}
\fancyhead[L]{\footnotesize 英语语法 名词性从句（订正）}
\fancyhead[R]{\footnotesize 赵文瀚}
\fancyfoot[C]{\thepage}
\renewcommand{\headrulewidth}{0.4pt}
\renewcommand{\footrulewidth}{0pt}

\begin{document}

\begin{center}
{\LARGE\bfseries 名词性从句}\\[0.2em]
{\large 错题订正与考点复盘}\\[0.3em]
{\footnotesize 2029 届高一英语 · 上海市大同中学 · 名词性从句专项（共 40 题）· 2026-10-28}
\end{center}

\vspace{-0.3em}
\noindent{\color{rulegray}\rule{\textwidth}{0.8pt}}

\vspace{0.35em}
\noindent\fcolorbox{black!15}{black!4}{%
  \parbox{\dimexpr\linewidth-2\fboxsep-2\fboxrule\relax}{%
    \small 日期：2026-10-28\qquad 来源：上海市大同中学高一英语「名词性从句」（作业号 260642，班级 6，共 40 题）\qquad 姓名：赵文瀚\\[0.25em]
    判错：\textcolor{wrongred}{\textbf{19}}\,/\enspace 40\qquad
    判对：\textbf{21}\,/\enspace 40\qquad
    （原卷右上角印 \textbf{53.8\%}；第 22 题一空对、一空错，按题计为判错）\qquad
    重做日期：\underline{\hspace{2.2cm}}}}

\vspace{0.4em}
\noindent{\footnotesize 说明：原卷已由老师红笔批改，下表「判定」\OK\ 判对、\NO\ 判错，\textbf{照录红笔}；\W{红色}＝被判错的作答，\A{加粗}＝正确答案。原卷\textbf{只打对错勾叉、未印标准答案}，「正确答案」为按标准语法补拟，第 16、19、26、27、30、40 题答案不唯一或与红笔口径有出入，见「原卷旁注」。}

\vspace{0.3em}
\renewcommand{\arraystretch}{0.98}
\noindent{\footnotesize\begin{tabular}{@{}L{0.60cm}@{\hspace{0.06cm}}L{8.30cm}@{\hspace{0.06cm}}L{2.20cm}@{\hspace{0.06cm}}L{2.30cm}@{\hspace{0.06cm}}L{0.85cm}@{\hspace{0.06cm}}L{2.85cm}@{}}
\toprule
\textbf{题号} & \textbf{题干（含空）} & \textbf{我的作答} & \textbf{正确答案} & \textbf{判定} & \textbf{错因} \\
\midrule
1  & \bk\ food is related to illness is not a new discovery. & That & That & \OK & — \\
2  & I had neither a raincoat nor an umbrella. That's \bk\ I got wet through. & why & why & \OK & — \\
3  & There are so many options available. I wonder \bk\ she will choose. & \W{what}（划改） & \A{which} & \NO & 有范围用 which \\
4  & The city is definitely no longer \bk\ it used to be. & what & what & \OK & — \\
5  & The suggestion \bk\ we go to visit the Natural Museum interested every one of us. & that & that & \OK & — \\
6  & It is still unknown \bk\ the plan will be carried out or not. & whether & whether & \OK & — \\
7  & \bk\ the workers insisted on was \bk\ they should be given more pay. & What；that & What；that & \OK & — \\
8  & Are you sure \bk\ she is honest? & that & that & \OK & — \\
9  & He made it quite clear \bk\ he wouldn't change his mind. & that & that & \OK & — \\
10 & ---What's that building? ---That's \bk\ the medical equipment is stored. & where & where & \OK & — \\
11 & It's none of your business \bk\ other people think about you? & \W{whether} & \A{what} & \NO & what 作介词宾语 \\
12 & We cannot figure out \bk\ quite a number of insects, birds, and animals are dying out. & \W{if} & \A{whether} & \NO & 正式文体不用 if \\
13 & Why not try your luck downtown, Bob? That's \bk\ the best jobs are. & \W{what} & \A{where} & \NO & 表地点用 where \\
14 & I have always been honest and straightforward, and it doesn't matter \bk\ I'm talking to. & \W{how}（前有涂改） & \A{who} & \NO & to 后缺宾语，指人 \\
15 & You may give the ticket to \bk\ you think is interested in the movie. & \W{if} & \A{whoever} & \NO & 缺主语、表无论谁 \\
16 & ---I went to see you at about ten this morning, but you were not in the office. ---Oh, that was probably \bk\ I was talking with the headmaster. & \W{that} & \A{because} & \NO & 表原因（when 亦可） \\
17 & So far there is no proof \bk\ people from other planets do exist. & that & that & \OK & — \\
18 & I was surprised by her words, which made me recognize \bk\ silly mistakes I had made. & what & what & \OK & — \\
19 & The difficulty we now meet with is \bk\ we can persuade him to tell the truth. & \W{whether} & \A{how} & \NO & 表方式，如何说服 \\
20 & \bk\ has a teaching qualification can apply for the job. & \W{Who} & \A{Whoever} & \NO & 主语从句表无论谁 \\
21 & \bk\ has won the first prize hasn't been announced yet. & Who & Who & \OK & — \\
22 & \bk\ we can't get seems better than \bk\ we already have. & What；\W{that} & What；\A{what} & \OK{}；\NO & 比较对象须用 what \\
23 & \bk\ nearly all behavior is learned behavior is a basic assumption of social scientists. & That & That & \OK & — \\
24 & Dr Black comes from either oxford or Cambridge, I can't remember \bk. & \W{whether} & \A{which} & \NO & 两者之一用 which \\
25 & I have several current magazines with me. You can take \bk\ you like most. & \W{what}（涂改） & \A{whatever} & \NO & 缺宾语，用 whatever \\
26 & \bk\ you can raise the audience's attention depends on the content of your speech. & \W{How much} & \A{Whether} & \NO & 主语从句表能否 \\
27 & The important thing is \bk\ name should be put on the top of the list. & \W{that} & \A{whose} & \NO & whose 作定语 \\
28 & It's generally considered unwise to give a child \bk\ he or she wants. & \W{what} & \A{whatever} & \NO & 泛指用 whatever \\
29 & \bk\ you study and work often means the difference between success and failure. & How & How & \OK & — \\
30 & As a new diplomat, he often thinks of \bk\ he can react appropriately on such occasions. & \W{how} & \A{whether} & \NO & 表是否（存疑） \\
31 & These photographs will show you \bk\ our village looks like. & \W{how} & \A{what} & \NO & look like 缺宾语 \\
32 & A story goes \bk\ Elizabeth I of England liked nothing more than being surrounded by clever and qualified noblemen at court. & that & that & \OK & — \\
33 & The other day, my brother drove his car down the street at \bk\ I thought was a dangerous speed. & what & what & \OK & — \\
34 & The question sometimes comes to my mind \bk\ it is worth the effort. & whether & whether & \OK & — \\
35 & Upon graduation he asked to be sent to \bk\ he is most needed. & where & where & \OK & — \\
36 & Perseverance is a kind of quality, and that's \bk\ it takes to do anything well. & what & what & \OK & — \\
37 & I told him that \bk\ I was surprised at was his attitude towards his study. & what & what & \OK & — \\
38 & The local government has promised to do \bk\ lies in its power to ease the hardships of the homeless. & \W{what} & \A{whatever} & \NO & 泛指用 whatever \\
39 & Mother made a promise \bk\ if I behaved well, she would buy me a new computer. & that & that & \OK & — \\
40 & \bk\ room this is should be ashamed of the mess it is in. & \W{Whose} & \A{Whoever's} & \NO & 表无论谁的 \\
\bottomrule
\end{tabular}}


\newpage
\vspace{0.28em}
\noindent{\footnotesize\textbf{原卷旁注}：（一）右上角印刷 \textbf{53.8\%} 为作业系统记录；按红笔逐题清点为 \OK\,21、\NO\,19（52.5\%；第 22 题一空对一空错，按题计判错），与 53.8\% 略有出入，请以老师讲评为准。（二）原卷\textbf{无标准答案栏}，红笔只打勾叉，故「正确答案」为补拟：第 16 题红笔判 \emph{that} 错，答案取 \emph{because}（\emph{when} 亦可）；第 26 题取 \emph{Whether}（\emph{How much} 若理解为"多少注意力"亦通，但 raise 不与 much 搭配）；第 27 题取 \emph{whose}（\emph{what name} 亦通）；第 30 题红笔判 \emph{how} 错，取 \emph{whether}（\emph{how} 作"如何"亦通，本条存疑）；第 40 题取 \emph{Whoever's}（\emph{Whose} 语法上亦成立，存疑）。（三）第 3 题先写 \emph{which} 后划掉改 \emph{what}，第 14、25 题作答有涂改，红笔均按最终作答判错。（四）原卷印题笔误照录：第 5 题 \emph{Natural Museum}、\emph{every one of us}，第 24 题 \emph{oxford} 小写、句间用逗号连接两分句。（五）第 7、22 题一句两空；第 7 题两空俱对，第 22 题第二空错。}

\begin{center}
{\Large\bfseries 错因归类与复盘}
\end{center}
\vspace{-0.45em}
\noindent{\color{rulegray}\rule{\textwidth}{0.8pt}}

\vspace{0.45em}
\noindent{\footnotesize 说明：按名词性从句考点归类 19 条错题（第 22 题计第二空一空），供汇总对账；逐题判定见前页。}

\vspace{0.4em}
\renewcommand{\arraystretch}{1.3}
\begin{center}
{\small\begin{tabular}{@{}p{6.4cm} p{8.2cm} c@{}}
\toprule
\textbf{易错考点} & \textbf{题号} & \textbf{题数} \\
\midrule
连接代词 what／whatever（含"缺宾语"判断） & 11, 25, 28, 31, 38 & 5 \\
连接代词 which／who／whoever／whose／Whoever's & 3, 14, 20, 24, 27, 40 & 6 \\
连接副词 where／how（表地点、方式） & 13, 19 & 2 \\
连接词 whether／if 之辨（介词后、主语从句、表"是否"） & 12, 26, 30 & 3 \\
其他连接词（because／that） & 16, 22 & 2 \\
\midrule
\textbf{合计} & \textbf{19 题} & \textbf{19} \\
\bottomrule
\end{tabular}}
\end{center}

\vspace{0.5em}
\noindent\textbf{关联知识点标签}
\vspace{0.2em}

\noindent\begin{tabular}{@{}p{8.2cm}@{\hspace{0.3cm}}p{8.2cm}@{}}
\begin{minipage}[t]{\linewidth}
\footnotesize
\STAR\ \textbf{看从句缺什么}：缺主语／宾语 $\to$ \textbf{what}（泛指、无先行词用 \textbf{whatever}）；缺定语且后接名词 $\to$ \textbf{whose}；不缺成分 $\to$ \textbf{that}（陈述）／\textbf{whether}（是否）\\[0.15em]
\STAR\ \textbf{有范围}：二者选一用 \textbf{which}；"无论谁"用 \textbf{whoever}／\textbf{whoever's}\\[0.15em]
\STAR\ \textbf{表语从句}：表地点 \textbf{where}、表原因 \textbf{because}、表方式 \textbf{how}
\end{minipage} &
\begin{minipage}[t]{\linewidth}
\footnotesize
\STAR\ \textbf{if 四不用}：介词后、主语从句、表语从句、与 or not 连用 $\to$ 一律 \textbf{whether}\\[0.15em]
\STAR\ \textbf{look like／think about 后缺宾语} $\to$ \textbf{what}，不用 how\\[0.15em]
\STAR\ \textbf{it doesn't matter}＋\textbf{who}（对象）；\textbf{It's none of your business}＋\textbf{what}
\end{minipage} \\
\end{tabular}

\vspace{0.6em}
\noindent\fcolorbox{black!15}{black!4}{%
  \parbox{\dimexpr\linewidth-2\fboxsep-2\fboxrule\relax}{%
    \small\textbf{复盘记录}\\[0.3em]
    最薄弱的板块：\underline{\hspace{9.0cm}}\\[0.5em]
    主要错误原因：\checkmarkbox\ 从句缺成分判断\quad \checkmarkbox\ what／whatever\quad \checkmarkbox\ whether／if\quad \checkmarkbox\ 连接副词\\[0.5em]
    本次复习的改进措施：\underline{\hspace{10.1cm}}\\[0.5em]
    \hspace{3.2em}\underline{\hspace{10.1cm}}\\[0.5em]
    下次复习日期：\underline{\hspace{2.2cm}}\qquad
    当前掌握度：\checkmarkbox\ 仍薄弱\quad \checkmarkbox\ 基本掌握\quad \checkmarkbox\ 已掌握
  }}

\end{document}
